\subsection{来源延迟下的有限签署容量}
\label{sec:eval-capital}

将组织 CAL 备付与授权 $B\in\{3600,28800\}$，
分别与 1 秒或 3 秒的来源公开延迟组合。
各成员拥有重叠额度 $\lfloor2B/3\rfloor$，即 2,400 或 19,200 CAL。
三个钱包执行相互独立的 A$\to$B$\to$C 单元，每单元本金 100 CAL，
两笔付款分别使用独立最终、用户出资的 FUEL 输入。
子交易使用父交易 TXCer 请求认证。父证书就绪后，独立计时器在指定延迟后公开父交易，
不依赖子交易进度。各组关闭自动来源补交以保持控制延迟，
不补资、不串行取证、不按时间解锁；Execution 与 Bytes 额度充足，未发生相应资源拒绝。

每格三轮全新运行，第二轮反转顺序。驱动在 60\,s 内按每秒 20 个单元
安排 1,200 个单元，未完成上限 128，窗口满时记录跳过的启动，最多排空 60\,s。
九个服务的委员会、成员、网关与 BlockStore 使用内存状态，钱包 NoSync；
本实验测量签署容量行为而非耐久存储延迟。

\begin{table}[t]\centering\small
\caption{每轮 1,200 个计划单元的有限资本结果。数值为三轮中位数 [最小值,最大值]。
部分批准按保留一份或两份成员批准的不同候选付款计数。}
\label{tab:capital-current}
\setlength{\tabcolsep}{3pt}
\begin{tabular}{@{}rrrr@{}}\toprule
$B$ / 延迟 & 启动 & 闭合 & 部分批准 \\\midrule
3,600 / 1 s & 212 [193,227] & 84 [65,99] & 63 [62,68] \\
3,600 / 3 s & 171 [166,208] & 43 [38,83] & 61 [55,64] \\
28,800 / 1 s & 1,200 & 1,200 & 0 \\
28,800 / 3 s & 1,200 & 1,200 & 0 \\
\bottomrule\end{tabular}\end{table}

$B=28800$ 时，六轮均闭合全部 1,200 个单元，没有 CAL 拒绝或残留部分批准。
1 秒延迟下的成员采样占用峰值为 5,000--5,200 CAL，3 秒时为 9,000--9,100 CAL，
均低于 19,200 CAL 额度（图~\ref{fig:capital-current}）。
$B=3600$ 时，各成员结束时均预留 2,400 CAL、可用为零。
停止状态审计核对了全部 96 份保留签名，分别落在一票或两票候选中；
每票 100 CAL 的预留恰好解释四名成员的重叠本地总量。
它们使用不同输入实例，没有成功输入冲突转换触发回收。
每轮 973--1,034 个计划单元未启动，125--128 个已启动单元在有界停止时仍未完成。

所有公开执行的付款仍然闭合，副本一致，最终公共缺口为零；
用户费用只扣一次，组织 FUEL 未被扣减。
因此，闭合恢复已完成批准的容量，其他部分批准仍保留占用。
高预算运行承接了该输入负载，低预算运行则表明覆盖安全本身不保证持续准入。
未完成请求保持删失记录，不赋予成功延迟。
补充材料~S4 保存逐成员预留、费用、失败及完整时间序列。

\begin{figure}[t]\centering
\includegraphics[width=\columnwidth]{capital-current-zh.pdf}
\caption{成员 CAL 预留除以额度：每秒内所有成员采样值的最大值。
每档预算包含三条运行轨迹（实线、虚线、点线）。60 秒准入窗口之后有界排空，
各轨迹在对应运行结束处停止。}
\label{fig:capital-current}
\end{figure}

